\chapter{Barmah Forest Virus}
\index{Barmah Forest Virus}

\section*{Definition and Overview}
Barmah Forest virus is a mosquito-borne virus found in Australia that causes a febrile illness with joint pain, rash, and fatigue. It belongs to the same family of viruses as Ross River virus and produces a very similar disease, known as Barmah Forest virus disease. The virus cannot be spread directly from person to person, and most people who become infected recover fully, although joint symptoms may persist for some time.

The condition is a notifiable disease in Australia, meaning cases are reported to public health authorities. There is no specific treatment or vaccine, so management focuses on relieving symptoms while the immune system clears the infection.

\section*{Epidemiology}
Barmah Forest virus occurs across Australia, particularly in the northern regions, the Murray-Darling basin, and parts of Western Australia and Queensland. Outbreaks are more common after heavy rainfall, which allows mosquito numbers to rise. People of any age and either sex can be infected, but disease is most often reported in adults who spend time outdoors in areas with mosquitoes. Each year several hundred to a few thousand cases are notified to health authorities, and the infection is a notifiable disease in all Australian states and territories. Infection is much more common than reported disease, since many cases cause few or no symptoms.

\section*{Aetiology and Pathophysiology}
\begin{itemize}
    \item \textbf{Cause:} Barmah Forest virus, an \textit{Alphavirus} transmitted by the bite of infected mosquitoes (mainly \textit{Culex} and \textit{Aedes} species)
    \item \textbf{Reservoir:} the virus is maintained in marsupial and other native animal hosts; humans are incidental hosts
    \item \textbf{Transmission:} mosquito bites, most often in late summer and autumn, or in the wet season in northern Australia
    \item \textbf{Mechanism:} the virus replicates at the bite site and in lymph nodes before spreading through the bloodstream; the immune response to the virus drives much of the joint pain and rash
\end{itemize}

\section*{Clinical Features}
\begin{itemize}
    \item Fever, chills, and headache
    \item Joint pain and swelling, particularly of the wrists, fingers, ankles, and knees
    \item Muscle aches and generalised fatigue
    \item A maculopapular rash (flat red spots and small bumps), often on the trunk, arms, and legs
    \item Tender, swollen lymph nodes
    \item Symptoms usually begin 3--11 days after the mosquito bite
\end{itemize}

\section*{Differential Diagnosis}
\begin{itemize}
    \item \textbf{Ross River virus disease:} clinically very similar and more common
    \item Other mosquito-borne infections, including dengue, Japanese encephalitis, and malaria in travellers
    \item Other causes of fever with rash and joint pain, such as rubella or parvovirus infection
    \item Bacterial infections such as leptospirosis or rickettsial disease (for example, Q fever or scrub typhus)
    \item Immune-mediated joint disease, such as early rheumatoid arthritis or reactive arthritis
\end{itemize}

\section*{When to Seek Medical Care}
See a doctor if you develop fever with joint pain, rash, or severe fatigue, especially if you have recently been bitten by mosquitoes or travelled within Australia.

\textbf{Contact your local emergency services for:}
\begin{itemize}
    \item High fever with severe headache and neck stiffness
    \item Confusion, drowsiness, or fitting
    \item Severe joint pain that prevents you from using a limb
    \item Difficulty breathing or chest pain
    \item Signs of dehydration from persistent vomiting
\end{itemize}

\section*{Diagnosis and Investigations}
\begin{itemize}
    \item \textbf{History and examination:} symptoms, mosquito exposure, travel within Australia, and examination of joints and skin
    \item \textbf{Blood tests:} serology to detect Barmah Forest virus antibodies (IgM and IgG) in the blood; a rising or high antibody level confirms recent infection
    \item \textbf{Serial testing:} a repeat blood test a few weeks later may be needed to confirm the diagnosis
    \item \textbf{Full blood count:} often normal or mildly abnormal, but helps exclude other infections
    \item \textbf{Ruling out other viruses:} testing may include Ross River virus, dengue, and other arboviruses because the diseases look alike
\end{itemize}

\section*{Management}
\begin{itemize}
    \item \textbf{Rest:} adequate rest, especially while feverish
    \item \textbf{Analgesics:} paracetamol for fever and pain; anti-inflammatory medicines such as ibuprofen may help joint pain, taken under medical advice
    \item \textbf{Fluids:} drink plenty of fluids to prevent dehydration
    \item \textbf{Joint care:} gentle movement and physiotherapy can help maintain joint function while symptoms settle
    \item \textbf{No antiviral treatment:} there is no specific antiviral medicine or vaccine; treatment is supportive
    \item \textbf{Follow-up:} review with a doctor if symptoms persist beyond a few weeks, particularly joint pain
\end{itemize}

\section*{Complications}
\begin{itemize}
    \item Persistent or recurrent joint pain and stiffness lasting weeks to months -- the most common complication
    \item Chronic fatigue that may take many months to resolve
    \item Secondary bacterial infection of scratched skin lesions (uncommon)
    \item Rarely, encephalitis (inflammation of the brain) or other neurological involvement
    \item Worsening of symptoms in people with pre-existing joint or immune conditions
\end{itemize}

\section*{Prognosis and Outcomes}
Most people recover fully within several weeks to a few months. Joint symptoms may persist or recur for months and occasionally for more than a year, but permanent joint damage is uncommon. Recovery is usually complete, and the infection provides lasting immunity to the virus. Outcomes are generally good, although some people find the prolonged fatigue and joint pain disruptive to daily activities.

\section*{Prevention and Self-Care}
\begin{itemize}
    \item Use insect repellent containing DEET, picaridin, or oil of lemon eucalyptus when outdoors
    \item Wear long sleeves, long trousers, and socks in mosquito areas, particularly at dawn and dusk
    \item Use bed nets or screens and insecticide sprays where mosquitoes are present
    \item Remove standing water around the home (buckets, tyres, plant saucers) where mosquitoes breed
    \item Limit outdoor activity when mosquito numbers are high after rain or flooding
    \item No vaccine is available, so personal protection against mosquito bites is the main preventive measure
\end{itemize}

\section*{Sources and Further Reading}
The following sources provide reliable, up-to-date information on Barmah Forest virus.
\begin{itemize}
    \item Centers for Disease Control and Prevention. \textit{Mosquito-borne diseases.} https://www.cdc.gov
    \item World Health Organization. \textit{Vector-borne diseases.} https://www.who.int
    \item MedlinePlus. \textit{Mosquito-borne diseases.} https://medlineplus.gov
    \item NHS. \textit{Infections.} https://www.nhs.uk/conditions/
\end{itemize}
